\section{Text Features}
\label{sec:text-features}
\slidesrc{Text features: 1--23; Blackboard: features}

Most NLP tasks that classify texts (by topic, by author, by sentiment, human or bot,
\ldots) need to turn a text into numbers first. Each number is a \textbf{feature}: a
measurable property of the text, such as the percentage of pronouns, the TF-IDF weight
of a word, the number of emotion words. The course map
(Section~\ref{sec:course-map}) listed the most common ones; this section explains them
and, above all, \emph{when} to use which.

\begin{nbdef}{Feature}
A \textbf{feature} of a text is a number (or a symbol) computed from the text that
describes one of its properties. A text is represented by a \textbf{feature vector}
$(f_1, f_2, \ldots, f_k)$, which is the input of a classifier or of another algorithm.
Some features belong to a \emph{single word} (being a function word), some to a
\emph{word in its sentence} (its POS), some to the \emph{whole document} (the
percentage of pronouns), and some need a whole \emph{corpus} (TF-IDF).
\end{nbdef}

This section covers:
\begin{enumerate}[nosep]
  \item the families of features: stylistic, content-based, content and style
        together (LIWC), vocabulary and idiosyncrasies, graph-based
        (Section~\ref{sec:feature-types});
  \item which features to use for which goal, with the summary table written on the
        blackboard (Section~\ref{sec:choosing-features}).
\end{enumerate}

\subsection{Families of Features}
\label{sec:feature-types}
\slidesrc{Text features: 2--13}

\subsubsection{Stylistic features}
\slidesrc{Text features: 2--8}

\textbf{Stylistic} features describe \emph{how} something is written, not \emph{what}
it is about: sentence length, punctuation, the proportion of pronouns or adjectives.
Two people writing about the same topic use different styles, and the same person
writing about different topics keeps much of their style. This is why stylistic
features are the basis of \textbf{author profiling} (guessing the age, gender, native
language, personality of the author) and \textbf{authorship attribution} (who wrote this
text?).

The slides show the 29 stylistic features used by an author profiling system for four
languages:

\begin{center}
\footnotesize
\begin{tabularx}{\linewidth}{@{}r X c | r X c@{}}
\toprule
\textbf{No.} & \textbf{feature} & \textbf{all 4} & \textbf{No.} & \textbf{feature} &
  \textbf{all 4} \\
\midrule
1 & \% question sentences & \checkmark & 16 & \% words with 3+ syllables & \\
2 & \% short sentences & \checkmark & 17 & average syllables per word & \\
3 & \% long sentences & \checkmark & 18 & \% adjectives & \\
4 & average sentence length & \checkmark & 19 & \% adverbs & \\
5 & average word length & \checkmark & 20 & \% capitals & \checkmark \\
6 & \% words with 6+ letters & \checkmark & 21 & \% colons & \checkmark \\
7 & \% words with 2--3 letters & \checkmark & 22 & \% determiners & \\
8 & \% semicolons & \checkmark & 23 & \% digits & \checkmark \\
9 & \% punctuation & \checkmark & 24 & \% full stops & \checkmark \\
10 & \% pronouns & & 25 & \% interjections & \\
11 & \% prepositions & & 26 & \% modals & \\
12 & \% coordinating conjunctions & & 27 & \% nouns & \\
13 & \% commas & \checkmark & 28 & \% personal pronouns & \\
14 & \% articles & & 29 & \% verbs & \\
15 & \% words with one syllable & & & & \\
\bottomrule
\end{tabularx}
\captionof{table}{The stylistic features of an author profiling system (from a paper
cited in the slides). All 29 are used for English; ``all 4'' marks those used also for
Dutch, Spanish and Italian. The others need a POS tagger or a syllable counter, which
the system had only for English.}
\label{tab:stylistic-features}
\end{center}

\begin{nbintuition}[why only some features in every language?]
Features 1--9, 13, 20, 21, 23, 24 can be computed by simply counting characters, words
and sentences: they work in any language. Features like ``\% adjectives'' need a POS
tagger (Section~\ref{sec:pos-tagging}) and ``syllables per word'' needs a syllable
counter, which are language-specific tools. A recurring lesson: the cheapest features
are the most portable.
\end{nbintuition}

\paragraph{N-grams and collocations.}
The co-occurrence of the same $n$ characters or words: character n-grams (\emph{ing\_},
\emph{\_th}) capture spelling and morphology habits, word n-grams and collocations
(\emph{strong tea}, Section~\ref{sec:corpora}) capture habitual phrasing.

\paragraph{POS (part of speech).}
Playing a particular role (``part'') in a sentence (``speech'') is a feature of
\emph{one word inside a sentence}: no need for the full document or corpus. The
\emph{number} of words with a given POS (number of pronouns, of adjectives, \ldots),
possibly normalised by the length of the document, is a feature of the
\emph{document}.

\paragraph{FW (function words).}

\begin{nbdef}{Function word}
In linguistics, \textbf{function words} (also called \emph{functors}) are words that
have little lexical meaning, or an ambiguous meaning, and express grammatical
relationships among other words within a sentence, or specify the attitude or mood of
the speaker.
\end{nbdef}

Being a function word is a feature of the \emph{word itself}: if the word is not
ambiguous, we do not even need the sentence it belongs to. The \emph{number} of function
words, possibly normalised by the document length, is a feature of the document. Kinds of
function words:
\begin{itemize}[nosep]
  \item \textbf{articles}: \emph{the}, \emph{a} (in some inflected languages the article
        takes the case of the following noun);
  \item \textbf{pronouns}: inflected in English (\emph{he}--\emph{him},
        \emph{she}--\emph{her});
  \item \textbf{conjunctions}: \emph{and} (uninflected in English);
  \item \textbf{subordinating conjunctions}: \emph{if}, \emph{then}, \emph{well},
        \emph{however}, \emph{thus};
  \item \textbf{auxiliary verbs}: part of the conjugation of main verbs, always
        inflected (\emph{has} gone, \emph{will} go);
  \item \textbf{particles}: \emph{up}, \emph{on}, \emph{down} (give \emph{up});
  \item \textbf{interjections}: sometimes called ``filled pauses'' (\emph{uh}, \emph{um});
  \item \textbf{expletives}: words that take the place of a sentence, among other
        functions;
  \item \textbf{pro-sentences}: \emph{yes}, \emph{okay}.
\end{itemize}
There is a \textbf{huge overlap with stop words} (Section~\ref{sec:stop-words-tfidf}):
the word cloud on the slide, of a novel, is dominated by \emph{the}, \emph{and},
\emph{to}, \emph{of}, \emph{a}, \emph{was}, \emph{her}, \emph{he}.

\begin{nbwarning}[stop words are not always noise]
For topic classification we \emph{remove} stop words because they say nothing about the
topic. For author profiling we \emph{count} them, because \emph{how often} an author uses
\emph{the}, \emph{of}, \emph{I}, \emph{she} is one of the best clues to who the author is.
The same words are noise for one task and signal for another.
\end{nbwarning}

\paragraph{Readability.}
\textbf{Readability} measures how easy a piece of text is to read. It can include
elements of complexity, familiarity, legibility and typography. Readability formulas
usually look at sentence length, syllable density and word familiarity; an example is
the Flesch--Kincaid grade level of Section~\ref{sec:course-map}. Online tools (e.g.\
\url{https://readable.com}) compute many such scores.

\subsubsection{Content-based features}
\slidesrc{Text features: 9}

\textbf{Content-based} features describe \emph{what} the text is about:
\begin{itemize}[nosep]
  \item the frequency of words belonging to some given, known \textbf{semantic category}
        (dictionary-based, topic-based): e.g.\ how many words are about food, sport,
        money;
  \item \textbf{emotion words} (based on Ekman's or Plutchik's basic emotions,
        Section~\ref{sec:course-map});
  \item \textbf{TF-IDF} (Section~\ref{sec:stop-words-tfidf}). Recall:
        \begin{align*}
          \text{tf-idf}(t, d, C) &= \text{tf}(t, d)\cdot\text{idf}(t, C),\\
          \text{tf}(t, d) &= \frac{\#\,t \text{ in } d}{\text{length}(d)}, \qquad
          \text{idf}(t, C) = \log\frac{\abs{C}}{\abs{\{d \in C : t \in d\}}};
        \end{align*}
  \item \textbf{LSA} (Latent Semantic Analysis, Section~\ref{sec:dsm-in-practice}):
        dense vectors of the words or of the document.
\end{itemize}

\begin{nbwarning}[LSA and word2vec]
The slide writes ``LSA, implemented by word2vec''. Strictly, they are two different
methods that produce the same kind of object, a dense vector per word: LSA reduces a
word--document matrix with SVD (Section~\ref{sec:dsm-in-practice}), word2vec trains a
small neural network to predict words from their contexts. Both are distributional, and
both can be used as content features.
\end{nbwarning}

\subsubsection{Content- and style-based features: LIWC}
\slidesrc{Text features: 10--11}

\textbf{LIWC} (Linguistic Inquiry and Word Count) is a tool with hand-built dictionaries
of psychologically meaningful word categories. It counts, for a text, the percentage of
words in each category: some categories are about style (I-words, function words),
others about content (social words, emotions, cognitive processes). It then combines
them into summary variables.

\begin{nbexample}{LIWC on a technical paragraph}
The slides feed LIWC a paragraph from a computer science paper (about the execution
time of some reasoners analysing log files). The result, compared with the average for
professional or scientific writing:
\begin{center}
\small
\begin{tabular}{@{}l rr@{}}
\toprule
\textbf{LIWC dimension} & \textbf{this text} & \textbf{average (scientific writing)} \\
\midrule
I-words (I, me, my)      & 0.0  & 0.63 \\
social words             & 0.0  & 7.62 \\
positive emotions        & 0.0  & 2.32 \\
negative emotions        & 1.1  & 1.45 \\
cognitive processes      & 17.0 & 7.52 \\
\midrule
\multicolumn{3}{@{}l}{\emph{summary variables}} \\
analytic                 & 99.0 & 92.57 \\
clout                    & 37.5 & 68.17 \\
authenticity             & 65.6 & 24.84 \\
emotional tone           & 11.9 & 43.61 \\
\bottomrule
\end{tabular}
\end{center}
Reading it: the paragraph is extremely \emph{analytic} (99), full of cognitive-process
words (\emph{because}, \emph{depends}, \emph{except}: 17\% vs.\ 7.5\%), impersonal (no
I-words, no social words) and emotionally flat (tone 11.9). This is what we expect from a
technical text, and LIWC finds it by counting words alone.
\end{nbexample}

\subsubsection{Vocabulary and idiosyncrasies}
\slidesrc{Text features: 12}

Features that capture personal or informal writing habits:
\begin{itemize}[nosep]
  \item \textbf{unique words}: the number of distinct words in the sample (vocabulary
        richness);
  \item \textbf{capital words}: the number of words (longer than 2) written only in
        capital letters (SHOUTING);
  \item \textbf{spelling errors}: the number of words not found in the dictionaries of
        the language;
  \item \textbf{character flooding}: \emph{sooooo saaaad\ldots};
  \item \textbf{neologisms}: newly invented words;
  \item \textbf{slang}: \emph{tanta roba!} (Italian slang, ``great stuff!'');
  \item \textbf{emoticons}: \texttt{>(}.
\end{itemize}

\subsubsection{Graph-based features}
\slidesrc{Text features: 13}

A text can also be turned into a \textbf{graph}, and the properties of the graph used as
features. The slide shows an example of \textbf{EmoGraph} for a Spanish sentence
(\emph{He estado tomando cursos en línea sobre temas valiosos que disfruto estudiando y
que podrían ayudarme a hablar en público}, ``I have been taking online courses about
valuable subjects that I enjoy studying and that might help me to speak in public''):
\begin{itemize}[nosep]
  \item the nodes are the \textbf{POS tags} of the words (detailed tags such as
        \texttt{NCMP000}, common noun masculine plural, or \texttt{VMIP1S0}, main verb
        indicative present first person singular) and \textbf{semantic and emotional
        labels} of the words (\emph{education}, \emph{language}, \emph{joy},
        \emph{positive}, \ldots);
  \item the edges follow the order of the words in the text: an edge from a tag to the
        next one, and from a word's tag to its semantic/emotional labels;
  \item graph measures (which nodes are most central, how connected the graph is, \ldots)
        become the features. They capture \emph{how} a person combines grammar and
        emotions, a stylistic fingerprint.
\end{itemize}

\begin{nbsummary}[families of features]
\begin{itemize}[nosep]
  \item Stylistic: n-grams/collocations, POS counts, function words (huge overlap with
        stop words), readability, sentence/word length, punctuation.
  \item Content-based: semantic categories, emotion words, TF-IDF, LSA.
  \item Content + style: LIWC (psychological word categories + summary variables).
  \item Vocabulary and idiosyncrasies: unique words, capitals, spelling errors,
        flooding, neologisms, slang, emoticons.
  \item Graph-based: e.g.\ EmoGraph (POS + emotion labels as a graph).
  \item Scope of a feature: the word (FW), the word in its sentence (POS), the document
        (counts), the corpus (TF-IDF).
\end{itemize}
\end{nbsummary}

\subsection{Which Features for Which Goal}
\label{sec:choosing-features}
\slidesrc{Text features: 14--23; Blackboard: features}

\subsubsection{Content, FW, POS: when, why?}
\slidesrc{Text features: 14--22}

``Well\ldots\ it depends on what your goal is.''

\paragraph{Topic classification: content matters.}
If you need to classify a document into categories such as ``music'', ``biology'',
``food'' (three examples of BabelNet categories), you probably do not care about function
words and POS: what matters are the content words. Two texts from the slides:
\begin{itemize}[nosep]
  \item \emph{``Baked sea bass with lemon caper dressing. To make the dressing, mix the
        oil with the lemon zest and juice, capers, mustard, \ldots''} (a recipe: food);
  \item \emph{``These Russian basses sing so low you'll think the world is ending.
        \ldots\ the sound of a Russian oktavist, singing a full octave below a normal
        bass\ldots''} (music).
\end{itemize}
Both contain \emph{bass}, which in BabelNet has several senses: \emph{bass, basso} (an
adult male singer with the lowest voice), \emph{sea bass} (the lean flesh of a saltwater
fish), \emph{freshwater bass} (a fish), \emph{bass voice} (the lowest adult male singing
voice). The content words around it (\emph{baked}, \emph{lemon}, \emph{capers} vs.\
\emph{sing}, \emph{octave}) both pick the right sense (word sense disambiguation,
Section~\ref{sec:word-sense-disambiguation}) and the right topic. Counting pronouns or
articles would not help at all.

\paragraph{Author profiling: function words and POS matter.}
But if you are profiling the \emph{author} of a document\ldots

\begin{nbfact}{Argamon, Koppel, Pennebaker and Schler (2009)}
``However, the preponderance of evidence suggests that the most consistently effective
features over a wide variety of authorship problems are \textbf{function words} and
individual \textbf{parts-of-speech}.''

\smallskip\noindent
{\footnotesize S.\ Argamon, M.\ Koppel, J.\,W.\ Pennebaker, J.\ Schler,
\emph{Automatically profiling the author of an anonymous text}, Communications of the
ACM 52(2), 2009.}
\end{nbfact}

\begin{nbexample}{female vs.\ male}
The slides show two abstracts of linguistics papers, one by a female and one by a male
author, about similar topics (stylistic analysis, rhetorical devices). First the
\textbf{personal pronouns} are highlighted (\emph{my}, \emph{I}, \emph{he},
\emph{their}), then the preposition \textbf{of}. The point: the topic of the two texts
is the same, so content features cannot tell the authors apart, but the \emph{rates} of
such small words can. In the large studies by Argamon, Koppel and colleagues, texts by
women use more personal pronouns, texts by men more determiners, numbers and
prepositions such as \emph{of} (words that specify nouns).
\end{nbexample}

\subsubsection{So, what?}
\slidesrc{Text features: 23}

\begin{itemize}[nosep]
  \item There is a possibly endless amount of features. Many are standard and well
        studied, and it is also clear which NLP tasks they suit best; others (EmoGraph,
        for example) are less standard.
  \item Anyone can use existing features, already implemented by existing tools.
  \item The real challenge is to \textbf{invent new features}, or well-calibrated
        combinations of existing ones, and to \textbf{demonstrate} that they are useful
        for some given NLP task, on some existing benchmark.
\end{itemize}

\subsubsection{The blackboard summary}
\slidesrc{Blackboard: features}

In class the lecturer summarised which features are useful for which tasks in a table
on the blackboard. On the board the features are the \emph{columns}; here the table is
transposed (one feature per row) so that it fits the page. Nothing else is changed.

\begin{nbblackboard}{Which feature is useful for which task}{Blackboard: features}
\small
\renewcommand{\arraystretch}{1.25}
\begin{tabularx}{\linewidth}{@{}>{\bfseries\raggedright\arraybackslash}p{3.0cm}
  >{\raggedright\arraybackslash}X >{\raggedright\arraybackslash}p{4.0cm}@{}}
\toprule
\textnormal{\textbf{feature}} & \textbf{useful for} & \textbf{useless or not
  applicable for} \\
\midrule
Function words & author profiling: $\bullet$ gender of [the author] / pron[ouns];
  $\bullet$ one specific author & topic classification or filtering \\
TF-IDF & topic classification & not applicable without a corpus; WSD \\
POS & author profiling; genre profiling; readability & {\color{nbgray}\itshape(empty)} \\
Period / stylistic features & Bob or bot?; human author profiling; document genre
  profiling & {\color{nbgray}\itshape(empty)} \\
Emotions & sentiment analysis; affective computing; legal cases; fake news &
  {\color{nbgray}\itshape(empty)} \\
Length of words & {\color{nbgray}\itshape(empty)} & {\color{nbgray}\itshape(empty)} \\
Variability of words & readability & {\color{nbgray}\itshape(empty)} \\
Slang, flooding, jargon & {\color{nbgray}\itshape(empty)} &
  {\color{nbgray}\itshape(empty)} \\
\bottomrule
\end{tabularx}

\smallskip
{\color{nbgray}\scriptsize On the board a double arrow links the headings ``Period /
stylistic features'' and ``Length of words'' (they belong together). In the first cell
the text after ``gender of'' is partly unreadable; brackets mark what is completed from
context.}
\end{nbblackboard}

\paragraph{Reading the blackboard, row by row.}
\begin{itemize}
  \item \textbf{Function words}. Useful for \textbf{author profiling}: to guess the
        \emph{gender} of the author (through the use of pronouns, as in the female vs.\
        male example) or to recognise \emph{one specific author} (authorship
        attribution: everybody has a personal rate of \emph{the}, \emph{of},
        \emph{however}). Useless for \textbf{topic classification or filtering}: function
        words are the same in every topic, which is why they are removed as stop words.
  \item \textbf{TF-IDF}. Useful for \textbf{topic classification}: the words with high
        TF-IDF in a document are the ones that characterise its topic. \textbf{Not
        applicable without a corpus}: the idf needs a collection of documents $C$; with a
        single text we cannot compute it. Useless for \textbf{WSD}: TF-IDF weighs a word
        for a whole document, it says nothing about which sense one occurrence has.
  \item \textbf{POS}. Useful for \textbf{author profiling} (Argamon et al.),
        \textbf{genre profiling} (a scientific paper has more nouns, a novel more verbs
        and pronouns, a recipe many imperatives) and \textbf{readability} (many nouns and
        long noun phrases make a text harder).
  \item \textbf{Period / stylistic features}. ``Period'' here means the sentence
        (Italian \emph{periodo}): sentence length, punctuation, question sentences, the
        stylistic features of Table~\ref{tab:stylistic-features}. Useful to tell
        \textbf{a human (Bob) from a bot}, for \textbf{human author profiling}, and for
        \textbf{document genre profiling}.
  \item \textbf{Emotions}. Emotion words are useful for \textbf{sentiment analysis}
        (positive or negative opinion), \textbf{affective computing} (systems that
        recognise and react to human emotions), \textbf{legal cases} (e.g.\ the emotional
        load of testimonies or of threatening messages) and \textbf{fake news} (fake news
        often uses a strongly emotional, alarming language).
  \item \textbf{Length of words}. No task written on the board; it is linked by the
        arrow to the stylistic features (word length is one of them, features 5--7 of the
        stylistic table) and enters readability formulas.
  \item \textbf{Variability of words} (vocabulary richness, the number of unique words).
        Useful for \textbf{readability}: a text that uses many different and rare words
        is harder to read.
  \item \textbf{Slang, flooding, jargon}: no task written on the board; from the
        previous subsection they are idiosyncrasies of informal writers (useful for
        author profiling, e.g.\ age, and for bot detection).
\end{itemize}

\begin{nbintuition}[the rule behind the table]
Ask: \emph{does the task depend on what is said, or on who says it and how?}
\begin{itemize}[nosep]
  \item \textbf{What} is said (topic, filtering): use \emph{content} features (TF-IDF,
        semantic categories); remove function words.
  \item \textbf{Who} says it and \textbf{how} (author, gender, human vs.\ bot, genre,
        readability): use \emph{style} features (function words, POS, sentence and word
        length, vocabulary richness, idiosyncrasies).
  \item \textbf{How the author feels} (sentiment, fake news, affective computing): use
        \emph{emotion} features.
\end{itemize}
\end{nbintuition}

\begin{nbsummary}[choosing features]
\begin{itemize}[nosep]
  \item Topic classification: content words (TF-IDF); function words and POS useless.
  \item Author profiling: function words and POS are the most consistently effective
        features (Argamon et al., 2009); female vs.\ male: pronouns vs.\ noun specifiers.
  \item TF-IDF needs a corpus; useless for WSD.
  \item Stylistic (sentence-level) features: bot detection, author and genre profiling.
  \item Emotions: sentiment analysis, affective computing, legal cases, fake news.
  \item Variability of words, POS: readability.
  \item The challenge: invent new features and prove they help on a benchmark.
\end{itemize}
\end{nbsummary}

